\documentclass[11pt]{article}
\usepackage[utf8]{inputenc}
\usepackage[T2A]{fontenc}
\usepackage[russian]{babel}
\usepackage{amsmath,amssymb}
\usepackage[margin=2.5cm]{geometry}
\usepackage{xcolor}
\definecolor{linkblue}{RGB}{30,80,160}
\usepackage[colorlinks=true, urlcolor=linkblue, linkcolor=black, citecolor=black, breaklinks=true]{hyperref}
\usepackage{booktabs}
\usepackage{enumitem}

\title{Отчёт о проделанной работе: privacy-preserving обучение флота роботов}
\author{Julia Tsvetkova}
\date{16 сентября 2026}

\begin{document}
\maketitle
\tableofcontents

\section{Контекст и цель}

Отправная точка — R\&D-заявка \emph{Privacy-Preserving Federated Conformal
Calibration} (MPC + Rényi DP) для Arenadata: калибровка неопределённости VLA-моделей
на edge-роботах через federated conformal prediction с секьюрной агрегацией
квантиля non-conformity score. Задача этого этапа работы — (1) оценить реальную
новизну заявки, (2) при необходимости перенаправить проект в более сильную научную
постановку, (3) построить исполнимый роадмап, (4) заложить математическую матчасть
под него, включая доработку MPC-слоя по итогам последнего созвона.

\section{Этап 1: аудит новизны исходной заявки}

Сверка с литературой ([arXiv:2302.06322] FedCP²-QQ, [arXiv:2306.05131] DP-FedCP,
оба 2023) показала: связка federated conformal prediction + differential privacy
уже опубликована с 2023 года. Федеративное обучение VLA — тоже закрыто (FedVLA,
ICCV 2025). Membership inference на VLA-политиках — закрыто в середине 2026
(две статьи: [arXiv:2605.07088], VLALeaks [arXiv:2606.15165]), причём с разрушительным
результатом: temporal-signal атака по действиям робота даёт AUC~0.999, а все
проверенные защиты (DP-SGD, шум, округление, MC dropout) либо не работают, либо
ломают точность управления.

\textbf{Вывод:} заявленный \guillemotleft research gap\guillemotright{} исходной
колоды не существует. Открытыми остаются: (C1) реконструкция обучающих данных из
политики (не просто обнаружение членства), (C2) систематизация того, почему FL/DP
это не останавливают, (C3) защита, снижающая утечку без потери точности управления
— именно это прямо названо нерешённой проблемой в VLALeaks.

\section{Этап 2: пивот и 9-месячный роадмап}

Сформулирован новый проект — \emph{Your Robot Remembers Your Home} — с целевыми
venue USENIX Security / IEEE S\&P (fallback: CCS, NDSS, NeurIPS D\&B, CoRL).
Роадмап на 13 недель + фаза масштабирования (мес. 4--9):

\begin{center}
\begin{tabular}{@{}p{2.1cm}p{3.3cm}p{8.5cm}@{}}
\toprule
Недели & Фаза & Итог \\
\midrule
1--2 & Воспроизвести и прочитать & OpenVLA-7B + LoRA на LIBERO, SR~$\ge$~70\%, one-pager
  модели угроз \\
3--6 & Воспроизвести атаки & Action-MSE, temporal-curvature, LiRA — таблица AUC,
  совпадающая с литературой \\
7--10 & Реконструкция и защиты & gradient matching (gray-box) + black-box
  реконструкция; targeted unlearning + латентный регуляризатор $\to$ Парето вне
  опубликованного фронтира \\
11--13 & Абляции и статья & 4-стр. workshop-paper подан — временная метка
  застолблена \\
\bottomrule
\end{tabular}
\end{center}

Разделение труда: Julia — ML/робототехника (обучение, атаки, защита, эксперименты);
научрук из security — формализация модели угроз, методология оценки
(TPR@low-FPR, shadow-модели), venue-стратегия.

\section{Этап 3: матчасть по федеративному обучению флота}

По запросу «проресёрчить федеративное обучение и как это применимо здесь»
подготовлено 5 документов (скилл \texttt{deep-research}, репозиторий
\texttt{D:\textbackslash research-agents}), каждый — обзор области + примеры +
тезаурус:

\begin{enumerate}[label=\textbf{0\arabic*.}, leftmargin=2.2cm]
  \item \textbf{FL-основы} — FedAvg [McMahan et al., 2017], non-IID-коррекции,
        связь с проектом (робот = клиент, что едет в облако — веса или скаляр
        калибровки, не сырые данные).
  \item \textbf{Архитектура флота} — три наблюдаемые схемы (fully local /
        edge-cloud hybrid / fleet-scale пул данных) с примерами (Gemini Robotics
        On-Device, Tesla FSD, Amazon DeepFleet, Open X-Embodiment/RT-X) и реальной
        рыночной оценкой Grand View Research: on-premise 73\% / cloud 27\% выручки
        рынка ИИ-в-робототехнике (2025) — с честной оговоркой, что это доля
        выручки, не доля архитектур, и таксономия не совпадает 1:1 с нашей.
  \item \textbf{Shamir / MPC} — секьюрная агрегация Bonawitz et al. 2017;
        установлено, что MPC-слой исходной заявки — не контрибуция, а прямое
        применение известного протокола; зафиксирован незакрытый gap —
        byzantine-роботы.
  \item \textbf{Методы агрегации} — FedAvg / secure / робастная (Krum, median) /
        иерархическая / асинхронная со staleness-взвешиванием; предложена
        client--edge--cloud топология вместо плоской звезды как более реалистичная
        для бытовых роботов.
  \item \textbf{Топологии связи} — звезда / иерархическая звезда / decentralized
        gossip / policy merging; следствие для модели угроз (gray-box vs black-box
        меняется в зависимости от топологии).
\end{enumerate}

\section{Этап 4: формализация MPC-слоя по итогам созвона}

По итогам последнего созвона (\guillemotleft хорошенько изучить математику
процесса и улучшить её\guillemotright) подготовлен документ~06
(\texttt{06\_prss\_he\_compartmented\_threshold.tex}), формализующий три пункта из
рабочих заметок:

\begin{itemize}
  \item \textbf{PRSS} (Cramer--Damgård--Ishai, TCC~2005, [arXiv через citation graph
        не индексируется — Springer LNCS~3378]) — механизм, которым $N$ участников
        генерируют свежие маски для суммирования \emph{без} сервера как участника:
        структура доступа и PRF-ключи определены только над множеством роботов,
        сервер не получает ключевого материала ни на каком этапе.
  \item \textbf{Threshold-HE} (threshold Paillier) — альтернатива маскированию:
        секретный ключ расшифровки сам разделён по Shamir, сервер складывает
        шифротексты по аддитивному гомоморфизму, расшифровка требует коалиции
        $t$ участников. Практическое преимущество для роботов с нерегулярным
        расписанием: синхронизация нужна только на этапе расшифровки, не на каждом
        раунде маскирования.
  \item \textbf{\guillemotleft Знают двое — знает семья\guillemotright} —
        формализовано через compartmented secret sharing (Simmons/Brickell,
        1990): введён \emph{домохозяйственный след} $h(S)$ — число независимых
        домохозяйств в скомпрометированном множестве $S$, а не сырое число роботов
        $|S|$, как корректная единица измерения силы коалиции.
\end{itemize}

\subsection*{Направление улучшения (arXiv MCP, 2023--2025)}

Целевой поиск через deep-research/arXiv MCP дал более сильную и современную
альтернативу compartmented-схеме 1990 года:

\begin{itemize}
  \item \textbf{Verifiable Weighted Secret Sharing} [arXiv:2505.24289, 2025] —
        вес $\omega_i=1/|H_{j(i)}|$ на робота делает правило
        \guillemotleft знают двое — знает семья\guillemotright{} прямым следствием
        весовой формулы, а не отдельной надстройкой; плюс верифицируемость против
        злонамеренного дилера и 100$\times$ выигрыш по коммуникации против наивной
        взвешенной VSS.
  \item \textbf{FedGT} [arXiv:2305.05506, 2023] — group testing поверх той же
        группировки по домохозяйствам закрывает byzantine-gap, отмеченный в
        документе~03, тем же организационным примитивом, без отдельного механизма.
  \item \textbf{Optimal Computational Secret Sharing} [arXiv:2502.02774, 2025] —
        граница эффективности по размеру доли $\frac{|S|+|K|}{t}$ как ориентир
        для реализации на ограниченном бортовом хранилище робота.
\end{itemize}

Предложенное совмещение этих трёх компонентов специально под флот роботов,
сгруппированных по домохозяйствам, в проверенной литературе не найдено — это
кандидат на самостоятельный вклад уровня C2/C3 внутри основного security-роадмапа,
не просто инженерная деталь.

\section{Статус на конец этапа}

\begin{center}
\begin{tabular}{@{}p{5.5cm}p{9.4cm}@{}}
\toprule
Документ & Статус \\
\midrule
Artifact-роадмап (HTML) & Опубликован, 9-месячный план по неделям \\
01 -- FL-основы & Готово, 0 ошибок компиляции \\
02 -- Архитектура флота & Готово, с примерами и рыночной оценкой \\
03 -- Shamir/MPC & Готово, byzantine-gap зафиксирован \\
04 -- Методы агрегации & Готово \\
05 -- Топологии связи & Готово \\
06 -- PRSS/HE/compartmented + улучшение & Готово, 6 стр., включает предложение
  на базе arXiv 2023--2025 \\
\bottomrule
\end{tabular}
\end{center}

\textbf{Дальше по плану:} недели 1--2 основного роадмапа (инфраструктура OpenVLA +
LIBERO, чтение) остаются следующим шагом; предложенная weighted-verifiable
конструкция из раздела~4 — кандидат на отдельную проверку внутри фазы~2 (недели
9--10, defense-эксперименты) до того, как она попадёт в финальный текст статьи.

\section{Все использованные источники}

Полный список, по этапам работы; источники, повторяющиеся в нескольких
документах-матчасти, приведены один раз.

\subsection*{Этап 1 --- аудит исходной заявки}

\begin{itemize}[leftmargin=1.3cm]
  \item Humbert et al. \emph{One-Shot Federated Conformal Prediction (FedCP\textsuperscript{2}-QQ)}.
        ICML 2023. \url{https://arxiv.org/abs/2302.06322}
  \item Plassier et al. \emph{Conformal Prediction for Federated Uncertainty
        Quantification Under Label Shift (DP-FedCP)}. ICML 2023.
        \url{https://arxiv.org/abs/2306.05131}
  \item \emph{FedVLA: Federated Vision-Language-Action Learning}. ICCV 2025.
        \url{https://arxiv.org/abs/2508.02190}
  \item \emph{ForgeVLA}. \url{https://arxiv.org/abs/2605.07474}
  \item \emph{Membership Inference Attacks on Vision-Language-Action Models}.
        \url{https://arxiv.org/abs/2605.07088}
  \item \emph{VLALeaks}. \url{https://arxiv.org/html/2606.15165v1}
  \item Carlini et al. \emph{Extracting Training Data from Large Language
        Models}. USENIX Security 2021.
        \url{https://www.usenix.org/system/files/sec21-carlini-extracting.pdf}
  \item \emph{SoK: Security and Privacy of Foundation-Model-Powered Robots}.
        \url{https://arxiv.org/pdf/2606.16788}
  \item \emph{Privacy Risks in Reinforcement Learning for Household Robots}.
        \url{https://arxiv.org/abs/2306.09273}
  \item \emph{Fine-Tuning Vision-Language-Action Models (OpenVLA-OFT)}.
        \url{https://arxiv.org/abs/2502.19645}
\end{itemize}

\subsection*{Этап 3 --- матчасть по FL (документы 01--05)}

\textbf{01 --- FL-основы}
\begin{itemize}[leftmargin=1.3cm]
  \item McMahan et al. \emph{Communication-Efficient Learning of Deep Networks
        from Decentralized Data (FedAvg)}. AISTATS 2017.
        \url{https://arxiv.org/abs/1602.05629}
  \item \emph{Momentum Benefits Non-IID Federated Learning}.
        \url{https://arxiv.org/abs/2306.16504}
  \item \emph{FedDRL}. \url{https://arxiv.org/abs/2208.02442}
  \item \emph{Client Adaptation improves FL with Simulated Non-IID Clients}.
        \url{https://arxiv.org/abs/2007.04806}
  \item \emph{Federated Learning Survey: Multi-Level Taxonomy of Aggregation
        Techniques}. \url{https://arxiv.org/abs/2511.22616}
\end{itemize}

\textbf{02 --- Архитектура флота}
\begin{itemize}[leftmargin=1.3cm]
  \item \emph{Learning While Deploying: Fleet-Scale RL for Generalist Robot
        Policies}. \url{https://arxiv.org/abs/2605.00416}
  \item \emph{Advancing Multi-Robot Networks via MLLM-Driven Sensing,
        Communication, and Computation}. \url{https://arxiv.org/pdf/2604.00061}
  \item \emph{Offload or Overload}. \url{https://arxiv.org/html/2603.18284}
  \item \emph{RAPID: Edge-Cloud Partitioned Inference for VLA}.
        \url{https://arxiv.org/pdf/2603.07949}
  \item \emph{Robot Fleet Learning via Policy Merging (FLEET-MERGE)}.
        \url{https://arxiv.org/abs/2310.01362}
  \item \emph{Open X-Embodiment: Robotic Learning Datasets and RT-X Models}.
        \url{https://arxiv.org/abs/2310.08864}
  \item Google DeepMind. \emph{Gemini Robotics On-Device}. deepmind.google
        (блог).
  \item \emph{Offline Edge AI Robotics: Air-Gapped Suitcase Robot on Jetson
        Orin NX}.
        \url{https://dasroot.net/posts/2026/05/offline-edge-ai-robotics-jetson-orin-nx-gemma-4-air-gapped/}
  \item \emph{Boosting Productivity with Industrial Robots}.
        \url{https://www.automate.org/news/industrial-robots-remote-monitoring-and-control-solution-boosts-productivity}
  \item \emph{Computer Vision at Tesla for Self-Driving Cars}. thinkautonomous.ai.
  \item \emph{How Amazon deploys robots in its operations facilities}.
        \url{https://www.aboutamazon.com/news/operations/how-amazon-deploys-robots-in-its-operations-facilities}
  \item Grand View Research. \emph{Artificial Intelligence In Robotics Market
        Report}.
        \url{https://www.grandviewresearch.com/industry-analysis/artificial-intelligence-ai-robotics-market-report}
\end{itemize}

\textbf{03 --- Shamir/MPC}
\begin{itemize}[leftmargin=1.3cm]
  \item A.~Shamir. \emph{How to Share a Secret}. CACM 22(11), 1979.
  \item Bonawitz et al. \emph{Practical Secure Aggregation for
        Privacy-Preserving Machine Learning}. ACM CCS 2017.
        \url{https://dl.acm.org/doi/10.1145/3133956.3133982}
  \item \emph{PQSF: Post-Quantum Secure Privacy-Preserving FL}. Scientific
        Reports, 2024.
  \item \emph{Group Verifiable Secure Aggregate FL Based on Secret Sharing}.
        \url{https://pmc.ncbi.nlm.nih.gov/articles/PMC11926380/}
  \item \emph{SoK: Secure Aggregation based on Cryptographic Schemes for FL}.
        EURECOM.
        \url{https://www.eurecom.fr/publication/6974/download/comsys-publi-6974.pdf}
  \item \emph{SABLE: Secure And Byzantine robust LEarning}.
        \url{https://arxiv.org/abs/2309.05395}
\end{itemize}

\textbf{04 --- Методы агрегации}
\begin{itemize}[leftmargin=1.3cm]
  \item \emph{Asynchronous Hierarchical Federated Learning}.
        \url{https://arxiv.org/abs/2206.00054}
  \item \emph{Timely Asynchronous Hierarchical FL: Age of Convergence}.
        \url{https://arxiv.org/abs/2306.12400}
  \item \emph{Asynchronous FL on Heterogeneous Devices: A Survey}.
        \url{https://arxiv.org/abs/2109.04269}
  \item \emph{Asynchronous Multi-Server FL for Geo-Distributed Clients}.
        \url{https://arxiv.org/pdf/2406.01439}
\end{itemize}

\textbf{05 --- Топологии связи}
\begin{itemize}[leftmargin=1.3cm]
  \item \emph{Topology-aware Federated Learning in Edge Computing: A
        Comprehensive Survey}. \url{https://arxiv.org/abs/2302.02573}
  \item \emph{Decentralized Federated Learning: A Survey and Perspective}.
        \url{https://arxiv.org/abs/2306.01603}
  \item \emph{Decentralized Federated Averaging via Random Walk}.
        \url{https://arxiv.org/abs/2508.21286}
  \item \emph{Graph-based Gossiping for Communication Efficiency in
        Decentralized FL}. \url{https://arxiv.org/abs/2506.10607}
  \item \emph{Privacy-preserving Decentralized FL over Time-varying
        Communication Graph}. \url{https://arxiv.org/abs/2210.00325}
  \item \emph{Secure Decentralized FL via Gossip and Virtual Voting}.
        \url{https://arxiv.org/abs/2607.08651}
\end{itemize}

\subsection*{Этап 4 --- формализация MPC (документ 06)}

\begin{itemize}[leftmargin=1.3cm]
  \item Cramer, Damgård, Ishai. \emph{Share Conversion, Pseudorandom
        Secret-Sharing and Applications to Secure Computation (PRSS)}. TCC 2005.
        \url{https://link.springer.com/chapter/10.1007/978-3-540-30576-7_19}
  \item \emph{Generalized Pseudorandom Secret Sharing and Efficient
        Straggler-Resilient Secure Computation}. Springer, 2021.
  \item P.~Paillier. \emph{Public-Key Cryptosystems Based on Composite Degree
        Residuosity Classes}. EUROCRYPT 1999.
  \item G.~Simmons. \emph{How to (Really) Share a Secret}. CRYPTO 1990;
        E.~Brickell. \emph{Some Ideal Secret Sharing Schemes}. EUROCRYPT
        1989/90.
  \item Shehata, Han, Thyagarajan. \emph{Verifiable Weighted Secret Sharing}.
        2025. \url{https://arxiv.org/abs/2505.24289}
  \item Xhemrishi et al. \emph{FedGT: Identification of Malicious Clients in FL
        with Secure Aggregation}. 2023. \url{https://arxiv.org/abs/2305.05506}
  \item Aureliano, Cohen, D'Oliveira. \emph{Optimal Computational Secret
        Sharing}. 2025. \url{https://arxiv.org/abs/2502.02774}
\end{itemize}

\end{document}
